% ==============================================================================
%                   Premium Mathematical Lecture Notes Template
% ==============================================================================
% Build: latexmk | Language: [english] or [czech] in \usepackage
% ==============================================================================

\documentclass[12pt, letterpaper, twoside]{book}

% Import all styling, math, and custom commands
% Set your language here: [english] or [czech]
\usepackage[english]{resources/style}

% ==============================================================================
\begin{document}
% ==============================================================================

% --- Title Page ---
\maketitlepage
{Lecture Notes}
{Subject Title}
{Author Name}
{\today}

% --- Content ---
\chapter{Introduction}
\section{Core Concepts}

\begin{objective}
  \begin{itemize}
    \item Understand the core principles of the Prism Clarity aesthetic.
    \item Learn to use the new minimalist table and figure environments.
  \end{itemize}
\end{objective}

\begin{overviewbox}
  Brief summary or overview of the current section.
\end{overviewbox}

\begin{termsbox}
  Keywords, important terms, concepts.
\end{termsbox}

\begin{definition}[Sample Definition]
  This is a sample definition using the custom \texttt{tcolorbox} environment.
\end{definition}
\sidenote{This is a sidenote placed elegantly in the wide margin, characteristic of Tufte textbooks. It can contain extra details, citations, or references.}

\begin{theorem}[Sample Theorem]
  This is a sample theorem.
\end{theorem}
\begin{intertextbox}
This is some normal text after the theorem.\sidenotemark
\end{intertextbox}
\sidenotetext{Here is a second sidenote on the same page. It should be marked with a dagger ($\dagger$) since it's the second note.}
\begin{intuition}
  Before jumping into the rigorous proof, here is the basic idea of why the theorem holds true.
\end{intuition}

\begin{proof}
  This is the rigorous proof of the theorem. We can also display a massive wide equation that breaks into the margin:
  
  \begin{wideequation}
    \int_{-\infty}^{\infty} e^{-x^2} dx = \sqrt{\pi} \quad \text{and} \quad \mathcal{F}\{f(t)\} = \int_{-\infty}^{\infty} f(t) e^{-i\omega t} dt \quad \text{across the full page!}
  \end{wideequation}
\end{proof}

\begin{algo}[Search Algorithm]
  \begin{algorithmic}[1]
    \Require A list of numbers $L$
    \Ensure The maximum number in $L$
    \State $max \gets L[0]$
    \For{\textbf{each} $x \in L$}
    \If{$x > max$}
    \State $max \gets x$
    \EndIf
    \EndFor
    \State \Return $max$
  \end{algorithmic}
\end{algo}

\begin{application}[Engineering]
  This theorem is heavily used in structural engineering to calculate load capacities.
\end{application}

\begin{styledmarginfigure}
  \centering
  \begin{tikzpicture}
    \draw[thick, accent, fill=accent!20] (0,0) circle (1.2cm);
    \node[accentdark, align=center] at (0,0) {\textbf{Tufte}\\\textbf{Margin}};
  \end{tikzpicture}
  \captionof{figure}{A margin figure perfectly placed and styled in the wide column.}
\end{styledmarginfigure}

\begin{marginequation}{Euler's Identity in the margin.}
  e^{i\pi} + 1 = 0
\end{marginequation}

\begin{history}
  First discovered in 1894, this concept revolutionized the field.
\end{history}

\begin{pitfall}[Common Error]
  Do not confuse this with the secondary principle, which only applies in a vacuum!
\end{pitfall}

\begin{note}
  This is an important note or remark.
\end{note}

\begin{keyreference}
  \fullcite{gottesman2001encoding}
\end{keyreference}

\begin{problem}[Sample Problem]
  Solve for $x$.
\end{problem}

\begin{solution}
  The solution goes here.
\end{solution}

\subsection{Minimalist Aesthetics}

\begin{intertextbox}
Here is an example of a minimalist table using alternating row shading and clean rules.
\end{intertextbox}
\begin{table}[htbp]
  \centering
  \rowcolors{2}{gray!10}{white}
  \begin{tabular}{llc}
    \toprule
    \textbf{Category} & \textbf{Environment} & \textbf{Color Group} \\
    \midrule
    Concepts          & Theorem, Definition  & Cobalt \& Cyan       \\
    Actions           & Problem, Pitfall     & Purple \& Coral      \\
    Context           & History, Reference   & Gold \& Indigo       \\
    \bottomrule
  \end{tabular}
  \caption{A clean, minimalist table without vertical rules.}
\end{table}

\begin{intertextbox}
Here is an example of a clean, minimalist geometric figure (Venn Diagram) using TikZ.
\end{intertextbox}
\begin{figure}[htbp]
  \centering
  \begin{tikzpicture}
    % Left circle
    \draw[thick, accent] (0,0) circle (1.5cm);
    \node[accentdark] at (-0.8, 0) {\textbf{Theory}};

    % Right circle
    \draw[thick, probcolor] (1.5,0) circle (1.5cm);
    \node[probcolor!80!black] at (2.3, 0) {\textbf{Practice}};

    % Intersection
    \node[font=\bfseries, text=titleblue] at (0.75, 0) {Notes};
  \end{tikzpicture}
  \caption{Venn diagram demonstrating the minimalist figure style.}
\end{figure}

\begin{summary}
  We have successfully introduced all 8 new environments, demonstrated the Prism Clarity aesthetic, and shown how Tufte-style margin notes and full-width blocks can be used!
\end{summary}

\section{Advanced Features}

\begin{theorem}[Nesting Example]
  The template now supports seamless nesting of environments via \texttt{nested\_parser.lua}.
  \begin{proof}
    This is a nested proof within a theorem. Notice how the left border elegantly combines the colors of both the parent theorem and the child proof.
  \end{proof}
\end{theorem}

\begin{exam}[Exam Strategy]
  \exlabel{Tip:} The \texttt{exam} environment is perfect for highlighting key takeaways or strategies for exams.
  \begin{examparts}
    \item Use the \verb|\exlabel| command for bold, colored inline headings.
    \item The \texttt{examparts} environment provides a clean, unbulleted list for structured advice.
  \end{examparts}
\end{exam}


\clearpage
% ==============================================================================
\end{document}
% ==============================================================================
